%% ma: L12-B08-0177
%% lop: 12
%% chuong: 2
%% bai: 8
%% dang: 12.8.53
%% loai: TN4
%% muc: VD
%% dap_an: "C"
%% nguon: "CĐ2.4.pdf giáo viên gửi 27/10/2026 (M07, Lê Minh Khôi – Nguyễn Anh Khôi, Chương II Chủ đề 4 Các bài toán về tâm tỉ cự), Câu 27"
%% kien_thuc: "Bài 8 (khoảng cách hai điểm, biểu thức tọa độ), Bài 6 (bình phương vô hướng, tâm tỉ cự), Bài 7 (hình chiếu lên mặt phẳng tọa độ)"
%% trang_thai: cho-duyet
%% ly_do: "Dạng mới 12.8.53 đề xuất, chờ giáo viên duyệt"
\begin{ex}
Trong không gian $Oxyz$, cho ba điểm $A(1;4;5)$, $B(0;3;1)$, $C(2;-1;0)$. Tọa độ của điểm $M$ thuộc mặt phẳng $(Oxz)$ để biểu thức $MA^2+2009MB^2+2027MC^2$ đạt giá trị nhỏ nhất là
\choice{$M\left(-\frac{4055}{4037};\frac{364}{367};\frac{2014}{4037}\right)$}{$M\left(\frac{4055}{4037};0;-\frac{2014}{4037}\right)$}{\True $M\left(\frac{4055}{4037};0;\frac{2014}{4037}\right)$}{$M\left(-\frac{4055}{4037};0;\frac{2014}{4037}\right)$}
\loigiai{Gọi $I=\dfrac{A+2009B+2027C}{4037}=\left(\dfrac{4055}{4037};\dfrac{4004}{4037};\dfrac{2014}{4037}\right)$.

Biểu thức bằng $4037MI^2+$ hằng số, nhỏ nhất khi $M$ là hình chiếu của $I$ lên $(Oxz)$: \[M\left(\dfrac{4055}{4037};0;\dfrac{2014}{4037}\right).\]}
\end{ex}
